% SPDX-FileCopyrightText: Copyright (c) 2024-2026 Yegor Bugayenko
% SPDX-License-Identifier: MIT

\section{Net}\label{sec:sockets}

Objects used for both server-side and client-side TCP sockets
  work in accordance with \citet{posix}.
All objects presented in this Section belong to the \ff{nk} package.

The \deff{socket} object is an abstraction of a TCP socket.
The \deff{connect} attribute establishes a connection.
It accepts an abstract object with one void attribute which is the socket descriptor.
The socket descriptor has the following attributes:

\begin{itemize}
    \item \deff{send}: send a message to the socket
    \item \deff{recv}: receive \(x\) amount of bytes from the socket
    \item \deff{as-input}: make ``input'' from socket
    \item \deff{as-output}: make ``output'' from socket
\end{itemize}

When you dataize \ff{socket.connect}, it creates a new socket,
  connects to the given IP on the given port, dataizes the provided scope,
  closes the socket and returns bytes retrieved from scope dataization.
An \adeff{error} is returned if any of the described steps failed.

This code opens a connection to TCP port 80 of \ff{localhost} and
  then reads the entire stream as a Unicode string:

\begin{ffcode}
(nk.socket "127.0.0.1" 80).connect
  io.tee-input > [s]
    s.as-input
    io.console
\end{ffcode}

The \deff{listen} attribute allows you to listen to incoming connections as a server.
It accepts an abstract object with one void attribute which is the socket descriptor.
The current socket descriptor has the same attributes as the descriptor
  from \ff{socket.connect} and also one extra attribute \deff{accept}.
This attribute is used to wait for an incoming connection.
It also accepts an abstract object with one void attribute
  which is the client socket descriptor, which was connected to your server.

This code emulates a server on TCP port 8080 of \ff{localhost},
  then it accepts an incoming connection and sends \ff{"Hello, world"} to it:

\begin{ffcode}
(nk.socket "127.0.0.1" 8080).listen
  s.accept > [s]
    client.send "Hello, world" > [client]
\end{ffcode}
